%% -*-latex-*-

\documentclass[a4paper,11pt]{article}

\usepackage[T1]{fontenc}
\usepackage[french]{babel}
\usepackage[utf8]{inputenc}
\usepackage{lmodern}
\usepackage{xspace}
\usepackage{xcolor}

\usepackage{mathpartir}
\usepackage{amssymb}
\usepackage{amsmath}

\usepackage{listings}
\usepackage{listingsutf8}

\lstdefinelanguage{OCaml}{
  morekeywords={
    and,as,assert,begin,class,constraint,do,done,downto,else,end,
    exception,external,false,for,fun,function,functor,if,in,
    include,inherited,initializer,lazy,let,match,method,module,
    mutable,new,object,of,open,or,private,rec,sig,struct,
    then,to,true,try,type,val,virtual,when,while,with
  },
  sensitive=true,
  morecomment=[s]{(*}{*)},
  morestring=[b]"
}

\lstset{
  language=OCaml,
  basicstyle=\small\ttfamily,
  keywordstyle=\sffamily\bfseries,
  commentstyle=\itshape,
  stringstyle=\ttfamily,
  columns=fullflexible,
  literate={é}{{\'e}}1,
  morekeywords={module}
}

\input{commands}
\input{ocaml_syntax}

\newcommand{\ibullet}{\item[$\bullet$]}
\newcommand{\inferlabel}[1]{\inferrule*[right=\textsc{#1}]}

\title{Correction du TP 1 d'Interprétation et compilation}
\author{Christian Rinderknecht}
\date{23 et 24 septembre 2003}

\begin{document}

\maketitle

\section{Questions}

\input{tp1-questions}

\section{Réponses}

\subsection{Module \textsf{Eval}}

Nous donnons ici le contenu du module \textsf{Eval} (fichier
\textsf{eval.ml}). Il contient plusieurs sous-modules correspondant
chacun à une calculette avec différen\-tes constructions qui sont
ajoutées au fur et à mesure des questions.

\lstinputlisting[
  language=OCaml,
  inputencoding=utf8
]{eval.ml}

\subsubsection{Sous-module \textsf{Eval.Initial}}
\label{Eval.Initial}

Le sous-module \textsf{Eval.Initial} regroupe les définitions de type
et de valeurs correspondant à la calculette avec constantes et
opérateurs arithmétiques. Il répond à la question~\ref{q1}.

L'exception \textsf{DivByZero} sera déclanchée lorsqu'une division par
zéro sera anticipée par la calculette.

Le dernier cas du filtre de la fonction \textsf{eval} est nécessaire car la
syntaxe abstraite, c'est-à-dire le type \textsf{Ast.expr}, inclut tous les
constructeurs pour réaliser l'ensemble du travail pratique.

\subsubsection{Sous-module \textsf{Eval.WithLet}}
\label{Eval.WithLet}

Le module \textsf{Eval.WithLet} reprend le module
\textsf{Eval.Initial} et l'étend avec les variables et la liaison
locale. Il répond à la question~\ref{q2}.

Les environnements $\rho$ sont codés de façon fonctionnelle.

La fonction \textsf{empty\_env} implante l'environnement vide.

La fonction \textsf{extend} implante l'extension des environnements
par une nouvelle liaison (il s'agit de l'opérateur $\oplus$). Ainsi
l'implantation de $\rho \oplus x \mapsto v$ est \textsf{extend env
  (x,v)}.

Il y a maintenant deux types d'erreurs: la division par zéro et les
variables libres (c'est-à-dire non introduites par un
\textsf{let}). Il est alors utile de regrouper ces erreurs sous un
même type. Ainsi, l'exception \textsf{Err (FreeVar x)} est déclanchée
quand la variable \textsf{x} est libre dans l'expression. Par exemple:
\texttt{let a = 1 in b} (ici la variable libre est \texttt{b}).

La fonction \textsf{eval} prend ici un argument supplémentaire:
l'environ\-nement dans lequel l'expression doit être évaluée.


\subsubsection{Sous-module \textsf{Eval.WithCondInt}}
\label{Eval.WithCondInt}

Le module \textsf{Eval.WithCondInt} reprend et étend le module
\textsf{Eval.WithLet} avec une construction conditionnelle
discriminant les branches avec un entier (à zéro ou non), dans le
style du langage~C. Il répond à la question~\ref{q3}. On nous donne:

\begin{itemize}

    \ibullet \textbf{Syntaxe concrète}
  {\small
   \begin{verbatim}
Expression ::= ...
  | "ifz" Expression "then" Expression "else" Expression
   \end{verbatim}
}

  \vspace*{-15pt}

  \ibullet \textbf{Exemple} \qquad \texttt{1+2*(ifz 0 then 3
  else 4)}

  \ibullet \textbf{Syntaxe abstraite}

  \kwd{type} \type{expr} \equal{} \texttt{...} \vbar{} \cst{Ifz}
  \kwd{of} \type{expr} \texttt{*} \type{expr} \texttt{*} \type{expr}

\end{itemize}

  \noindent Une sémantique opérationnelle pour notre conditionnelle est:

  \begin{mathpar}
  \inferlabel{if-then}
    {\eval{\rho}{e_1}{\cst{Int} \, (n)}\\
     n \neq 0\\
     \eval{\rho}{e_2}{v_2}}
    {\eval{\rho}{\cst{Ifz} \, (e_1,e_2,e_3)}{v_2}}
  \and
  \inferlabel{if-else}
    {\eval{\rho}{e_1}{\cst{Int} \, (0)}\\
     \eval{\rho}{e_3}{v_3}}
    {\eval{\rho}{\cst{Ifz} \, (e_1,e_2,e_3)}{v_3}}
  \end{mathpar}

L'implantation par un nouveau cas dans le filtre de la fonction
\textsf{eval} est immédiat.


\subsubsection{Sous-module \textsf{Eval.WithCond}}
\label{Eval.WithCond}

Le module \textsf{Eval.WithCond} reprend le module
\textsf{Eval.WithCondInt} et l'étend avec une conditionnelle qui
discrimine ses branches à l'aide d'une expression booléenne (comme en
O'Caml, par exemple). Il répond à la question~\ref{q4}. Nous devons
ajouter au langage de la calculette les expressions booléennes.  Pour
faciliter le travail, cela est déja fait au niveau de la syntaxe
concrète et abstraite -- cf. \textsf{lexer.mll}, \textsf{parser.mly}
et \textsf{ast.ml}. Dans ce dernier, la syntaxe abstraite contient la
définition des booléens.

Voici les syntaxes reconnues par défaut par la calculette:
  \begin{itemize}

    \ibullet \textbf{Syntaxe concrète}
  {\small
   \begin{verbatim}
Boolean ::= true | false
      | Boolean "and" Boolean
      | Boolean "or" Boolean
      | "not" "(" Boolean ")"
Expression ::=  ...
      | "if" Boolean "then" Expression "else" Expression
   \end{verbatim}
}

  \vspace*{-6pt}

  \ibullet \textbf{Exemple} \qquad \texttt{1+2*(if true then 3
  else 4)}

  \ibullet \textbf{Syntaxe abstraite}

  \begin{tabbing}
    \kwd{type} \type{boolean} \= \equal{} \= \cst{True} \vbar{}
    \cst{False} \vbar{} \cst{Not} \kwd{of} \type{boolean}\\
    \> \vbar \> \cst{And} \kwd{of} \type{boolean} \(\times\)
    \type{boolean}\\
    \> \vbar \> \cst{Or} \kwd{of} \type{boolean}
    \(\times\) \type{boolean}\\
    \kwd{type} \type{expr} \equal{} \texttt{...} \vbar{} \cst{If}
    \kwd{of} \type{boolean} \(\times\) \type{expr} \(\times\)
    \type{expr}
  \end{tabbing}

  \end{itemize}

On remarque que la solution présentée ici ne permet pas qu'un booléen
soit une expression, donc le type des valeurs restera inchangé dans la
sémantique, mais il faudra une seconde relation d'évaluation dédiée
aux booléens. Donc, en théorie, il faudrait choisir une autre notation
que $\eval{\rho}{e}{v}$, mais en pratique c'est commode de reprendre
la même, à ceci près que les expressions booléennes n'ont pas besoin
d'un environnement pour s'évaluer car nous n'avons pas permis au
\texttt{let ... in ...} de lier des booléens. Donc nous noterons la
relation d'évaluation des booléens $\evalb{b}{\overline{b}}$, où $b$
est une valeur O'Caml de type \textsf{Ast.bexpr} et $\overline{b}$ est de
type \type{bool}.

  \begin{mathpar}
  \inferlabel{if-then}
    {\evalb{b}{\kwd{true}}\\
     \eval{\rho}{e_2}{v_2}}
    {\eval{\rho}{\cst{If} \, (b,e_2,e_3)}{v_2}}
  \and
  \inferlabel{if-else}
    {\evalb{b}{\kwd{false}}\\
     \eval{\rho}{e_3}{v_3}}
    {\eval{\rho}{\cst{If} \, (b,e_2,e_3)}{v_3}}
  \end{mathpar}

L'implantation de cette sémantique est directe (elle s'appuie
inductivement sur la syntaxe abstraite). l'implantation par un nouveau
cas dans le filtre de la fonction \textsf{eval} est immédiat.

\subsubsection{Sous-module \textsf{Eval.WithEnvList}}
\label{Eval.WithEnvList}

Le module \textsf{Eval.WithEnvList} reprend le module
\textsf{Eval.WithCond} en remplaçant le codage fonctionnel de
l'environnement par un codage à l'aide d'une structure de donnée
nommée \emph{liste d'association}. Il répond à la
question~\ref{q5}. Une liste d'association est une liste de paires qui
modélisent les liaisons. Ainsi, $\rho(x)$ est implanté non plus par
\textsf{env x} (\textsf{env} est une fonction OCaml) mais par
\textsf{List.assoc x env} (\textsf{env} est une liste
d'association). La raison est essentiellement l'efficacité de la
recherche d'une liaison, bien qu'il serait encore plus efficace
d'utiliser une autre structure de donnée, comme une table de hachage,
si le nombre de variable est élevé en moyenne par programme.

Les environnements $\rho$ sont codés à l'aide d'une liste
d'association. L'envi\-ron\-nement vide est alors simplement codé par
\textsf{[\;]} (plus besoin de \textsf{empty\_env}).

Au lieu de coder $\rho(x)$ par \textsf{env x} (plus un
\textbf{\textsf{try with}}, on fait \textsf{lookup x env}.

L'ajout d'une nouvelle liaison devient l'ajout d'une paire an tête de
liste: $\rho \oplus x \mapsto v$ est codé par \textsf{(x,v)::env} au
lieu de \textsf{extend env (x,v)} avec un codage fonctionnel. On n'a
donc plus besoin de la fonction \textsf{extend}.

Dans la définition de la fonction \textsf{eval}, le motif \textsf{Var
  x} change et appelle \textsf{lookup x env} dorénavant, et le cas de
la liaison locale est plus simple.

\subsection{Module \textsf{Main}}

Le module \textsf{Main} est le pilote de la calculette. C'est ce
module \emph{uniquement} qui interagit avec l'usager et, en
particulier, filtre les erreurs des autres modules (ici
\textsf{Eval}). Il est composé de plusieurs sous-modules correspondant
chacun à un des sous-modules de \textsf{Eval} de même nom. Par exemple
\textsf{Main.Initial} correspond à (en l'occurrence \emph{utilise})
\textsf{Eval.Initial}.

\lstinputlisting[
  language=OCaml,
  inputencoding=utf8
]{main.ml}

\subsubsection{Sous-module \textsf{Main.Initial}}
\label{Main.Initial}

Le sous-module \textsf{Main.Initial} correspond au cas initial avec
les quatre opérations arithmétiques de base ainsi que les
constantes. Il répond donc à la question~\ref{q1} en complétant le
module \textsf{Eval.Initial} (cf. section~\ref{Eval.Initial}). La
seule erreur qui puisse provenir de \textsf{Eval.Initial} est
\textsf{DivByZero}.

\subsubsection{Sous-module \textsf{Main.WithLet}}
\label{Main.WithLet}

Le sous-module \textsf{Main.WithLet} correspond au module
\textsf{Eval.WithLet} (cf. section~\ref{Eval.WithLet}). En complétant
ce dernier, il répond ainsi à la question~\ref{q2}. Les variables et
la liaison locale sont possibles, donc il faut filtrer en plus le cas
des variables libres.

\subsubsection{Sous-module \textsf{Main.WithCondInt}}
\label{Main.WithCondInt}

Le sous-module \textsf{Main.WithCondInt} correspond à
\textsf{Eval.WithCondInt} (cf. section~\ref{Eval.WithCondInt}). En
complétant ce dernier, il répond ainsi à la question~\ref{q3}. Du
point de vue du pilotage, rien ne change: les erreurs possibles sont
toujours les mêmes qu'avec \textsf{Main.WithLet}.

\subsubsection{Sous-module \textsf{Main.WithCond}}\label{Main.WithCond}

Le sous-module \textsf{Main.WithCond} correspond à
\textsf{Eval.WithCond} (cf. section~\ref{Eval.WithCond}). En
complétant ce dernier, il répond ainsi à la question~\ref{q4}.

\subsubsection{Sous-module \textsf{Main.WithEnvList}}
\label{Main.WithEnvList}

Le sous-module \textsf{Main.WithEnvList} correspond à
\textsf{Eval.WithEnvList} (cf. section~\ref{Eval.WithEnvList}). En
complétant ce dernier, il répond ainsi à la question~\ref{q5}. Nous
devons néanmoins réécrire \textsf{Main.WithEnvList.calc} car le codage
de l'environ\-ne\-ment vide à changé (c'est \textsf{[\;]} maintenant).

La fonction \textsf{main} est celle qui est appelée après le lancement
du programme. Pour tester en boîte blanche (c.-à-d. en modifiant le
code) il faut alors modifier les appels à \textsf{calc} et
recompiler. Par exemple, nous avons ici l'appel à
\textsf{WithEnvList.calc}, que l'on peut remplacer par un appel à
\textsf{WithCond.calc} etc. Cette méthode n'est pas bonne car, en
général, un test en boîte grise ou noire est préférable. Pour cela il
faudrait proposer,par exemple, à l'usager des options sur la ligne de
commande de lancement de la calculette.

Cette liaison globale est le point d'entrée du programme (il s'agit
d'évaluer l'appel à \textsf{main}). Son intérêt est de demander à
l'environnement d'exécution de tracer toutes les exceptions qui
pourraient remonter de l'appel à \textsf{main}.


\end{document}
